\section{Doubly Linked List: Construction and Operations}
\label{sec:ll-doubly-basics}

\problemstatement{Implement a doubly linked list (DLL): build it, insert
before or after a node, and delete a node -- all in $O(1)$ once the node is
in hand, thanks to the backward \texttt{prev} pointer.}

\begin{patternbox}
A DLL adds a \texttt{prev} pointer to each node. The pay-off is $O(1)$
deletion of a known node (no need to walk to find its predecessor) and easy
backward traversal -- which is exactly why LRU caches and browser-history
structures use it.
\end{patternbox}

\subsection*{Two pointers per node}

Each node stores \texttt{prev} and \texttt{next}. Inserting a node $x$
between $a$ and $b$ rewires four pointers: $a.\texttt{next}=x$,
$x.\texttt{prev}=a$, $x.\texttt{next}=b$, $b.\texttt{prev}=x$. Deleting a
node $x$ splices its neighbours together: $x.\texttt{prev}.\texttt{next}=
x.\texttt{next}$ and $x.\texttt{next}.\texttt{prev}=x.\texttt{prev}$, with
null checks at the ends. Because a node knows its own predecessor, deletion
needs only the node itself.

\begin{ideabox}[The prev Pointer Buys O(1) Deletion]
In a singly linked list, deleting a node requires its predecessor, found by
an $O(n)$ walk. The \texttt{prev} pointer stores that predecessor, so any
node can be removed in constant time -- the feature that makes DLLs the
backbone of caches.
\end{ideabox}

\subsection*{C++ implementation}

\begin{lstlisting}
#include <bits/stdc++.h>
using namespace std;

struct Node {
    int val;
    Node* prev;
    Node* next;
    Node(int v) : val(v), prev(nullptr), next(nullptr) {}
};

Node* build(const vector<int>& a) {
    Node dummy(0), *tail = &dummy;
    for (int x : a) {
        Node* node = new Node(x);
        tail->next = node; node->prev = tail;
        tail = node;
    }
    Node* head = dummy.next;
    if (head) head->prev = nullptr;
    return head;
}

Node* deleteNode(Node* head, Node* x) {           // O(1) unlink
    if (x->prev) x->prev->next = x->next;
    else head = x->next;                          // x was the head
    if (x->next) x->next->prev = x->prev;
    return head;
}

void insertAfter(Node* a, int v) {                // splice x between a and a->next
    Node* x = new Node(v);
    x->next = a->next; x->prev = a;
    if (a->next) a->next->prev = x;
    a->next = x;
}

void print(Node* head) {
    for (Node* c = head; c; c = c->next) cout << c->val << " ";
    cout << "\n";
}

int main() {
    Node* head = build({1, 2, 3});
    insertAfter(head, 9);                          // 1 9 2 3
    print(head);
    head = deleteNode(head, head->next);           // remove the 9
    print(head);                                   // 1 2 3
    return 0;
}
\end{lstlisting}

\begin{complexitybox}
\textbf{Time Complexity.} $O(1)$ per insert/delete of a known node; $O(n)$
to build. \textbf{Space Complexity.} $O(1)$ extra per operation.
\end{complexitybox}

\begin{takeawaybox}
A doubly linked list is a singly linked list with a backward pointer,
trading a little memory for $O(1)$ deletion and two-way traversal. Keep the
four-pointer insert and two-pointer delete rewiring precise -- off-by-one
null handling at the ends is the usual bug.
\end{takeawaybox}
